\scanpage{78}
Il s'ensuit que le quotient de
$\unsure{\mathcal X}$ par antipodisme est le $=$ que
précédemment, et que les suites de
\unsure{p.r.|p.v.} intervenant dans un circuit
sont les \inserted{mêmes} $\unsure{\Rightarrow|=}$ pour les deux --- donc
notre recette n'est pas plus satisfaisante,
en ce qui concerne les circuits
superréguliers ! Il faudrait modifier
3), en prenant la prescription opposée
pour définir le transport d'orientation,
pour trouver \deleted{\unsure{de}} bons circuits
superréguliers.

Qu'en serait-il dans le cas
où $\uline{\Sigma}$ est un syst. de \uline{droites} dans un
vrai plan projectif réel $X$ --- quand on regarde
le système de droites dans $X^\vee$ correspondant
aux sommets de $\Sigma$, i.e. les droites joignant
deux à deux les pts $\delta_i$ correspondant aux
droites $D_i$ ? On suppose \inserted{au besoin} que les $\delta_i$ sont non
alignés $3$ à $3$ (i.e. $\uline{\Sigma}\ \unsure{AdM}$), et que
les sommets \inserted{distincts des $\delta_i$} (de ce syst. de
$\frac{n(n-1)}2$ droites) sont simples (quitte
à perturber un peu le syst. des
$\delta_i$). Donc il n'y a que des
sommets d'ordre $4$, en dehors des
sommets ``\unsure{super-singuliers|superréguliers}'' $\delta_i$.

Considérons une face $F$ de $\uline{\Sigma}^\vee$,
qui définit donc une \unsure{p.r.|p.v.}
$D=D_F$, et considérons les
arêtes incidentes \inserted{$a_\alpha$ en $\overline F_\alpha=\Delta_\alpha\cap\overline F$} qui correspondent
à \uline{certaines} spécialisations
$D\to D_\alpha$ (où $\alpha$ est l'indice d'un
sommet de $\Sigma$, $\Delta_\alpha$ la droite
corr. dans $X^\vee$). Le sommet
$s_\alpha$ est le sommet d'un
tri. de spécialisation pour $D$,
soit $a'_\alpha$ l'arc de spécialisation,
correspondant. Je voudrais
\margin{77}
